% ============================================================
% PRÁCTICAS INTEGRADORAS: IMPORTACIÓN POR COURIER Y COMPARACIÓN DE PRÉSTAMOS
% ============================================================
\clearpage
\section*{PRÁCTICAS INTEGRADORAS}
\subsection*{Comercio exterior y finanzas con datos reales}

\begin{recuadro}
Estas dos prácticas aplican lo visto en las unidades de Sistema Impositivo, Comercio Exterior y
Préstamos a situaciones reales de un técnico que equipa su propio taller. Se trabajan con
precios, cotizaciones y tasas \textbf{del día}, buscados por cada estudiante, y se resuelven en
una planilla de cálculo.
\end{recuadro}

% ------------------------------------------------------------
\section*{PRÁCTICA 1: IMPORTAR UN COMPONENTE POR COURIER}

\subsubsection*{Situación}
Para reparar un equipo de comunicaciones VHF, tu taller necesita un módulo que no se consigue
en el país. Lo encontrás en una tienda de Estados Unidos y decidís traerlo por el
\textbf{régimen de courier} (pequeños envíos).

\subsubsection*{Ejemplo resuelto}
\begin{center}
\small
\begin{tabular}{|p{0.60\linewidth}|r|}
\hline
\textbf{Concepto} & \textbf{Importe} \\ \hline
Precio del módulo en la tienda (FOB) & USD 250,00 \\ \hline
Envío (flete) hasta Argentina & USD 60,00 \\ \hline
\textbf{Valor en aduana (CIF = producto + flete + seguro)} & \textbf{USD 310,00} \\ \hline
¿Supera la franquicia de USD 400? & No \\ \hline
IVA 21\% sobre el valor CIF & USD 65,10 \\ \hline
\textbf{Total en dólares} (producto + envío + IVA) & \textbf{USD 375,10} \\ \hline
Cotización del dólar del día (dato de ejemplo) & \$1.450 \\ \hline
\textbf{Costo total en pesos} & \textbf{\$543.895} \\ \hline
\end{tabular}
\end{center}

\paragraph*{Conclusiones del ejemplo.}
\begin{itemize}
  \item El precio de la tienda (USD 250) \textbf{no es} el costo real: con envío e impuestos el
        módulo costó un \textbf{50\% más} (USD 375,10).
  \item Si el envío hubiera superado los USD 400, sobre el excedente se pagan además derechos
        de importación (consultar la tasa vigente en ARCA).
  \item Se usaron 1 de los 5 envíos anuales con franquicia del titular.
  \item Si la tienda cobra con tarjeta en dólares, revisar además las percepciones o el tipo de
        cambio que aplica el banco.
  \item El courier puede cobrar gastos de gestión aduanera adicionales: hay que preguntarlo antes.
\end{itemize}

\subsubsection*{Consigna}
\begin{enumerate}
  \item Elegir un componente o instrumento real para el taller (un módulo, un conector
        especial, un analizador, una estación de soldado) en una tienda del exterior.
  \item Anotar precio, costo de envío y plazo de entrega.
  \item Calcular el valor CIF, el IVA y, si corresponde, los derechos sobre el excedente.
  \item Convertir a pesos con la cotización del día (indicar fuente y fecha).
  \item Buscar el mismo producto (o uno equivalente) en el país y comparar precio y plazo.
  \item Elegir el Incoterm que corresponde a la compra (¿quién paga el flete? ¿dónde se
        transfiere el riesgo?).
  \item \textbf{Conclusión:} ¿conviene importarlo o comprarlo en el país? Justificar con
        números y considerar también el plazo y la garantía.
\end{enumerate}

% ------------------------------------------------------------
\section*{PRÁCTICA 2: COMPARAR PRÉSTAMOS PARA EQUIPAR EL TALLER}

\subsubsection*{Situación}
Necesitás \textbf{\$3.000.000} para comprar un osciloscopio y un banco de pruebas. Tenés tres
opciones: dos préstamos personales a 12 meses y una compra en cuotas con tarjeta.

\subsubsection*{Ejemplo resuelto (sistema francés, 12 cuotas)}
\begin{center}
\small
\begin{tabular}{|p{0.30\linewidth}|r|r|r|}
\hline
 & \textbf{Banco A} & \textbf{Banco B} & \textbf{Tarjeta C} \\ \hline
Monto & \$3.000.000 & \$3.000.000 & \$3.000.000 \\ \hline
TNA & 45\% & 55\% & ``sin interés'' \\ \hline
TEA & 55,5\% & 71,2\% & --- \\ \hline
Cuota mensual & \$315.037 & \$330.569 & \$270.000 \\ \hline
Cantidad de cuotas & 12 & 12 & 12 \\ \hline
\textbf{Total a pagar} & \textbf{\$3.780.443} & \textbf{\$3.966.826} & \textbf{\$3.240.000} \\ \hline
Intereses (total $-$ monto) & \$780.443 & \$966.826 & \$240.000 (*) \\ \hline
\end{tabular}
\end{center}
{\footnotesize (*) En la opción C el negocio ofrece precio de contado \$3.000.000 y en 12 cuotas
``sin interés'' \$3.240.000: el recargo del 8\% sobre el precio de contado \textbf{es un
interés}, aunque no se llame así.}

\paragraph*{Conclusiones del ejemplo.}
\begin{itemize}
  \item Entre A y B, una diferencia de 10 puntos en la TNA significa \textbf{\$186.383 más}
        a pagar en un año.
  \item La \textbf{TEA} es siempre mayor que la TNA porque los intereses se capitalizan mes a
        mes; para comparar préstamos se usa la TEA, y mejor aún el \textbf{CFT}, que suma
        seguros, comisiones e impuestos.
  \item Las ``cuotas sin interés'' no siempre lo son: hay que comparar con el \textbf{precio de
        contado}.
  \item Antes de decidir, comprobar que la cuota entre en el presupuesto mensual del taller:
        se recomienda que no supere el 30\% de los ingresos.
\end{itemize}

\subsubsection*{Consigna}
\begin{enumerate}
  \item Entrar a los \textbf{simuladores de préstamos} de tres entidades (por ejemplo, un banco
        público, un banco privado y una billetera virtual) y simular \$3.000.000 a 12 meses.
  \item Registrar para cada una: TNA, TEA, \textbf{CFT}, cuota inicial, total a pagar y
        requisitos (ingresos mínimos, antigüedad, tipo de cuenta).
  \item Armar en la planilla la tabla de amortización de una de ellas por sistema francés y
        por sistema alemán, y comparar el total de intereses.
  \item Buscar una oferta de ``cuotas sin interés'' para el mismo equipo y calcular el interés
        oculto comparando con el precio de contado.
  \item \textbf{Conclusión:} ¿qué opción conviene? ¿Conviene endeudarse o ahorrar unos meses y
        comprar de contado? Justificar con números.
\end{enumerate}
